\section{Project Identity}
\label{sec:doctrine-identity}

\subsection{Purpose}
\label{sec:doctrine-identity-purpose}

The \texttt{doctrine/} corpus within \texttt{\textasciitilde/process-intelligence} is the immutable
research authority layer for the full-lifecycle process intelligence program. Its purpose is
singular and total: to define, in canonical and unrebased form, the foundational law stack
that every downstream repository must treat as authoritative. No downstream refactor mandate
may be issued, no M\&A claim manufactured, and no board-level projection emitted until this
layer speaks first.

The corpus currently comprises 34 canonical markdown files covering every dimension of the
law stack: the Blue River Dam upstream closure law, the MAPE-K autonomic knowledge actuation
loop, the twelve-stage process lifecycle calculus, the Evidence chain from raw records to
board-admissible receipts, the five process maturity levels, the algorithm taxonomy spanning
all fifteen process mining algorithm families, and the formal object taxonomy governing typed
execution authority. Each file carries explicit paper citations grounding every claim in
published process mining literature, including van der Aalst (2016, 2023), Carmona et al.\
(2018), Adriansyah et al.\ (2011), Leemans et al.\ (2013--2022), Berti et al.\ (2019--2023),
and Pesic et al.\ (2006).

The immutability covenant is structural: doctrine files may only be appended with dated
addendums and must never be rebased. This is not a preference; it is a manufacturing
precondition. An audit trail that can be rewritten is not an audit trail.

\subsection{Repository Surface}
\label{sec:doctrine-identity-surface}

The doctrine corpus occupies the \texttt{doctrine/} directory of the
\texttt{\textasciitilde/process-intelligence} research foundry. The 25 source files examined
for this chapter include: \texttt{BLUE\_RIVER\_DAM.md},
\texttt{AUTONOMIC\_KNOWLEDGE\_ACTUATION.md}, \texttt{CONFORMANCE\_AS\_LAW.md},
\texttt{EVIDENCE\_CHAIN.md}, \texttt{VAN\_DER\_AALST\_CANON.md},
\texttt{OBJECT\_CENTRIC\_SUPREMACY.md}, \texttt{MA\_READY\_PROCESS\_INTELLIGENCE.md},
\texttt{RECEIPT\_DOCTRINE.md}, \texttt{DOWNSTREAM\_AUTHORIZATION\_LAW.md},
\texttt{ALGORITHM\_TAXONOMY.md}, \texttt{PROCESS\_INTELLIGENCE\_DEFINED.md},
\texttt{GRADUATION\_LAW.md}, \texttt{PROCESS\_INTELLIGENCE\_SPR\_THESIS.md},
\texttt{PROCESS\_LIFECYCLE\_ONTOLOGY.md}, \texttt{PROCESS\_TRUTH\_AUTHORITY.md},
\texttt{FORMAL\_OBJECTS\_TAXONOMY.md}, \texttt{NAMED\_LAW\_REFUSAL.md},
\texttt{wasm-boundary-law.md}, \texttt{RESEARCH\_AUTHORITY.md},
\texttt{full-lifecycle-process.md}, \texttt{spr\_thesis\_actuation.md},
\texttt{lattice-monotonicity-verification.md}, \texttt{lifecycle\_algorithms.md},
\texttt{bpmn-or-join-completion.md}, and the sealed checkpoint
\texttt{checkpoints/PROCESS\_INTELLIGENCE\_ALIVE\_001.md}.

The repository surface is deliberately narrow: no source code, no compiled binaries, no
integration tests. It is a doctrine surface only. Its authority derives from the quality and
sourcing of its claims, not from execution artifacts. The execution authority belongs to
\texttt{wasm4pm}; the type foundry belongs to \texttt{wasm4pm-compat}; the manufacturing
machinery belongs to \texttt{ggen}. The doctrine corpus authorizes all three but executes
none of them.

\subsection{Primary Research Role}
\label{sec:doctrine-identity-role}

The thesis role of the doctrine corpus is \emph{immutable-process-law-foundation}. It
occupies the apex of the research authority stack, as documented in
\texttt{RESEARCH\_AUTHORITY.md}:

\begin{quote}
\texttt{\textasciitilde/process-intelligence} issues verdicts, authorizes downstream change,
and constitutes the research authority. Every downstream repository derives its authority
from doctrine and must not exceed that authority without explicit research program
authorization.
\end{quote}

In the dissertation's argument, the doctrine corpus serves three functions. First, it
establishes the formal vocabulary --- the precise type-theoretic and process-mining-theoretic
meanings of terms such as \emph{admitted evidence}, \emph{named law refusal},
\emph{graduation bridge}, and \emph{receipt-bearing execution}. Second, it provides the
falsifiability surface: every doctrine claim is either sourced to a published paper or
explicitly marked as a gap awaiting evidence. Third, it constitutes the checkpoint surface:
the \texttt{PROCESS\_INTELLIGENCE\_ALIVE\_001} checkpoint, sealed at 588 commits with 12 of
12 gate criteria verified, is the boundary condition that authorized the six downstream
workflows described in \texttt{DOWNSTREAM\_AUTHORIZATION\_LAW.md}.

\subsection{Key Artifacts}
\label{sec:doctrine-identity-artifacts}

The key artifacts of the doctrine corpus are organized by function:

\paragraph{Foundational Law Files.} \texttt{BLUE\_RIVER\_DAM.md} establishes the upstream
closure law and the five maturity levels. \texttt{EVIDENCE\_CHAIN.md} defines the typed
one-way lifecycle: $\text{Raw} \to \text{Parsed} \to \text{Admitted} \to
\{\text{Projected} \mid \text{Exportable} \mid \text{Receipted}\} \mid \text{Refused}$.
\texttt{CONFORMANCE\_AS\_LAW.md} frames conformance checking as process-level law
enforcement, not diagnostics.

\paragraph{Operating Equation.} The Blue River Dam operating equation, expressed across
multiple doctrine files as $\kappa(\rho(\alpha(\mu(O^*)))) \to \text{ALIVE} \mid
\text{PARTIAL} \mid \text{REFUSED}$, is the central formal claim: manufacture from knowledge
corpus, actuate, emit evidence, gate the outcome --- with no silent success permitted.

\paragraph{Checkpoint.} \texttt{PROCESS\_INTELLIGENCE\_ALIVE\_001.md} is the sealed phase
checkpoint carrying verification codes for all twelve gate criteria, including admissibility
boundary validation, autonomic actuation typestate verification, token-game fitness
confirmation, OCPQ refinement, and decommissioning retirement closure.

\paragraph{Authorization Law.} \texttt{DOWNSTREAM\_AUTHORIZATION\_LAW.md} is the active
legal surface that translates the sealed checkpoint into six specific authorized downstream
workflows, each with named targets, gate conditions, and receipt requirements.
